% Generated draft: requires mathematical and source review.
\documentclass[11pt]{article}
\usepackage[margin=25mm]{geometry}
\usepackage{fontspec}
\setmainfont{Latin Modern Roman}
\usepackage{amsmath,amssymb,mathtools,unicode-math}
\setmathfont{Latin Modern Math}
\usepackage{enumitem,needspace,xurl,hyperref}
\setlength{\parindent}{0pt}
\setlength{\parskip}{5pt}
\setlength{\emergencystretch}{3em}
\newcommand{\N}{\mathbb{N}}
\newcommand{\Z}{\mathbb{Z}}
\newcommand{\Q}{\mathbb{Q}}
\newcommand{\R}{\mathbb{R}}
\newcommand{\C}{\mathbb{C}}
\newcommand{\F}{\mathbb{F}}
\newcommand{\A}{\mathbb{A}}
\newcommand{\PP}{\mathbb{P}}
\newcommand{\E}{\mathbb{E}}
\newcommand{\eps}{\varepsilon}
\newcommand{\dd}{\,\mathrm{d}}
\begin{document}
\section*{UnsolvedMath ID: 20002452}
\textbf{Source problem code:} catalogue entry
\subsection{Definitions and mathematical statement}
The opening ``By what mechanism does this happen and how can one capture the general dynamics?'' has no antecedent in the canonical excerpt. The aid identifies it as the end of Kapitula's preceding vortex-pattern-return question, and the following H. Warchall entry as unrelated. The reported recovery from the preceding item is not actually included here, so the original Gross--Pitaevskii equation, sign convention, parameters, observed pattern, and notion of return remain unspecified.

The aid gives conditional diagnostics rather than a mechanism established for a particular PDE. Let $\Phi_t$ be a unique reversible flow on a complex phase space, equivariant under a group $G$ and gauge multiplication by $U(1)$. Define projective isotropy by
\[
G_{[u]}=\{g\in G:gu=e^{i\chi}u\text{ for some }\chi\in\R\}.
\]
Then $G_{[\Phi_tu]}=G_{[u]}$ while the flow and its inverse exist. Thus literal exact projective symmetry cannot disappear and return under these hypotheses. Forward equivariance alone preserves any initial exact symmetry; equality additionally uses reversibility. Approximate or density-based visual symmetry is a different observable.

A proposed reversible beating diagnostic is the Hermitian two-mode system
\[
i\dot a=\begin{pmatrix}\delta&\kappa\\\kappa&-\delta\end{pmatrix}a,
\quad a(0)=(1,0),\quad\delta,\kappa\in\R,\quad\kappa\ne0,\quad
\Omega=\sqrt{\delta^2+\kappa^2}.
\]
It predicts
\[
|a_2(t)|^2=\frac{\kappa^2}{\Omega^2}\sin^2(\Omega t),\qquad
T=\frac\pi\Omega,\qquad D_{\max}=\frac{\kappa^2}{\Omega^2},\qquad a(T)=-a(0).
\]
The return is up to gauge, and this is a normal form, not a derived reduction of the missing PDE. If unresolved modes produce forcing $r(t)$ in $i\dot a=Ha+r(t)$ with the same Hermitian $H$, the reported Duhamel control is
\[
\|a(t)-e^{-itH}a(0)\|\leq\int_0^t\|r(s)\|\,\dd s.
\]
A symmetry-equivariant linear coupling cannot connect inequivalent irreducible exact-symmetry representations; such a two-mode model requires compatible modes or explicitly broken symmetry. The aid recommends charge-preserving comparisons of fields, gauge/rotation quotients, vortex configurations modulo permutation, and densities, and an energy-norm test separating unitary beating from nonnormal Hamiltonian amplification. Its precise three diagnostic distances and energy norm are not supplied, so density resemblance must not be called exact field return.

The separate H. Warchall setup is retained. For $x\in\R^N$, $u(x,t)\in\R^M$, it gives
\[
Ju_t=\Delta u+g(u)\quad\text{(NLS)},\qquad
u_{tt}=\Delta u+g(u)\quad\text{(NLKG)},
\]
where $J$ is invertible skew-symmetric and $g(y)=h(|y|^2)\widehat y$, $\widehat y=y/|y|$ for $y\ne0$, is continuous. A value at zero compatible with continuity is required. The standing-wave ansatz is
\[
u(x,t)=e^{\mu Kt}\widehat\psi(x/|x|)w(|x|),\qquad
K=J^{-1}\text{ for NLS},\quad K^2=-I\text{ for NLKG},
\]
with $K$ real skew-symmetric, $\mu\in\R$, and $\widehat\psi:S^{N-1}\to S^{M-1}$ satisfying
\[
\Delta_S\widehat\psi=-\ell(\ell+N-2)\widehat\psi.
\]
The allowed $\ell$ are restricted by $M$, as the source notes. The printed profile equation is
\[
w''+\frac{N-1}{r}w'-\frac{\ell(\ell+N-2)}{r^2}w+f(w)=0,
\qquad f(w)=g_{\mathrm{sc}}(w)+\omega w,
\]
with printed $\omega=\mu$ for NLS and $\omega=\mu^2$ for NLKG. Here $g_{\mathrm{sc}}$ denotes the scalar radial restriction of the vector nonlinearity, resolving an overload in the extraction. There is an unresolved sign mismatch: the stated NLS equation and $K=J^{-1}$ directly give $g_{\mathrm{sc}}(w)-\mu w$, whereas the printed rule gives a plus. The sign is flagged rather than silently changed. The NLKG plus-$\mu^2$ rule is compatible with $K^2=-I$.

The source describes Galilean or Lorentz boosts as generating travelling waves, gives the motivating NLS form $u=e^{i\omega t}e^{im\vartheta}w(r)$ for $-iu_t-\Delta u=g(u)$, and reports smooth exponentially localized profiles under appropriate hypotheses. It lists $f'(0)<0$ and $F(t)=\int_0^t f(s)\,\dd s>0$ for some $t>0$ as essential conditions, without a complete existence theorem, origin conditions, or growth assumptions. These profile facts do not supply the missing vortex-return dynamics.
\subsection{Short English statement}
The opening vortex-return question is missing its setup. The aid offers symmetry and two-mode diagnostics, not a derived PDE mechanism. A separate encapsulated-vortex standing-wave presentation is preserved with its references and scalar-profile conditions, including an unresolved NLS sign mismatch.
\subsection{Sources}
\begin{itemize}
\item Local UnsolvedMath registry, permanent ID $20002452$, \texttt{catalogue entry}; \texttt{aim-pdes-notes.json}, \texttt{notes[63]}, source item 16.
\item American Institute of Mathematics, \emph{Stability criteria for multi-dimensional waves and patterns}, \url{https://aimath.org/WWN/multidimwaves/multidimwaves.pdf}. Supplied verification reports comparison of page 7 and preceding \texttt{AIM-PDES-0063}, identifying a final Kapitula question followed by H. Warchall's ``Stability of Encapsulated-Vortex Solutions.'' The promised recovered Kapitula entry is absent.
\item J. Iaia and H. A. Warchall, \emph{Encapsulated-vortex solutions to equivariant wave equations: existence}, \emph{SIAM J. Math. Anal.} 30 (1999), 118--139.
\item R. L. Pego and H. A. Warchall, \emph{Spectrally stable encapsulated vortices for nonlinear Schr\"odinger equations}, \emph{J. Nonlinear Sci.} 12 (2002), 347--394.
\item Research attempt 1: partial result, problem partially solved at run, rigor label ``proved,'' novelty confidence low. Registry classification partial progress, status partially solved, statement status reconstructed/unverified. No authoritative separate citation for the diagnostic theorem or observed return is supplied.
\end{itemize}
\end{document}
